\begin{abstract}
\addtocontents{toc}{\vspace{1em}}

\begin{spacing}{1.02}
Anti-Money Laundering (AML) in cryptocurrency networks presents a critical challenge at the crossroads of financial compliance, graph theory, and deep learning. Most existing machine learning methods treat transaction networks as fixed graph structures, which does not reflect the ongoing flow of funds and temporal sequencing characteristic of complex laundering operations. Moreover, conventional models tend to yield only a single binary illicit/licit classification that is not typologically differentiated and lacks the interpretable evidentiary subgraphs needed for actionable regulatory compliance and Suspicious Activity Report (SAR) filing.

To overcome such weaknesses, this dissertation introduces \textbf{TemporalAML}, an end-to-end framework unifying continuous temporal representation learning, multi-typology detection, and forensic explainability. The framework constructs a directed, weighted temporal graph from the benchmark Elliptic Bitcoin Dataset, which consists of 203,769 transaction nodes, 234,355 fund-flow edges, and 49 chronological time steps---incorporating 172 node attributes and six engineered topological features under a leak-proof, chronological assessment approach. At its architectural heart, TemporalAML presents a novel \textit{Learnable Fourier Time Encoding} mechanism where frequency parameters $\boldsymbol{\omega}$ and phase shifts $\boldsymbol{\phi}$ are trained directly within a multi-head Temporal Graph Attention Network (TGAT). By optimizing these temporal parameters via backpropagation, the model adaptively adjusts to the non-stationary, bursty velocity profiles typical of cryptocurrency transactions. The output of a shared latent representation is passed to parallel classification heads specialized for three different laundering typologies: \textit{Circular Transfers}, \textit{Layering Chains}, and \textit{Smurfing/Fan-Out} structures, utilizing class-weighted multi-task training to address extreme class imbalance. Model interpretability is achieved through an adapted GNNExplainer module, which extracts compact, time-ordered evidentiary subgraphs and automatically creates standardized SAR narratives for flagged alerts.

In Phase~1 of this dissertation, the theoretical foundation is established alongside the complete architectural implementation and empirical validation of the proposed methodology. Preliminary experiments show that the learnable Fourier time encoding successfully captures temporal dependency signals with positive gradient flow, outperforming fixed-frequency baselines on chronological splits. The resulting framework provides a robust basis for comprehensive benchmark evaluations, hyperparameter ablation studies, and regulatory usability assessments to be carried out during Dissertation Phase~2.

\vspace{0.5em}
\noindent\textbf{\textit{Keywords}:} Anti-Money Laundering, Temporal Graph Attention Networks, Learnable Fourier Time Encoding, Graph Neural Networks, Cryptocurrency, Bitcoin, GNNExplainer, Multi-Pattern Detection, Elliptic Bitcoin Dataset, Suspicious Activity Reports.
\end{spacing}

\end{abstract}